\documentclass[aspectratio=169]{beamer}
\usetheme{metropolis}
\usepackage{appendixnumberbeamer}
\usepackage{booktabs, hyperref}

\input{mgmt675-style}

\usepackage{tikz}
\usetikzlibrary{shapes.geometric, arrows.meta, positioning, calc}

% Title info
\subtitle{MGMT 675: Generative AI for Finance}
\title{Module 5: CLAUDE.md and Skills}
\author{Kerry Back}
\date{}

\begin{document}

\maketitle

\begin{frame}{Why Skills Matter for Finance}
When you build DCF models and portfolio optimizations in ad-hoc conversations, what happens when you start a new conversation?
\vspace{0.3cm}
\begin{columns}[T]
\begin{column}{0.5\textwidth}
\begin{shadedbox}[title=\textbf{The Problem}]
\begin{itemize}\small
  \item One conversation uses 10\% WACC, another uses 8\%
  \item One assumes 3\% terminal growth, another 2.5\%
  \item Portfolio optimization: 5-year history vs.\ 3-year
\end{itemize}
\end{shadedbox}
\end{column}
\begin{column}{0.5\textwidth}
\begin{shadedbox}[title=\textbf{The Fix: Skills}]
\begin{itemize}\small
  \item A DCF skill specifies: which assumptions, how to estimate them, what output format
  \item A portfolio skill specifies: data source, estimation window, constraints
  \item \alert{Every conversation follows the same playbook}
\end{itemize}
\end{shadedbox}
\end{column}
\end{columns}
\vspace{0.3cm}
\alert{No audit trail, no reproducibility --- unless you encode the methodology in a skill.}
\end{frame}

\begin{frame}{Finance Skill Examples}
\begin{columns}[T]
\begin{column}{0.5\textwidth}
\begin{shadedbox}[title=\textbf{DCF Skill}]
\begin{itemize}\small
  \item Sales growth from trailing 5-year CAGR
  \item Sector median EV/EBITDA for exit multiple
  \item Two-way sensitivity table (WACC vs.\ growth)
\end{itemize}
\end{shadedbox}
\end{column}
\begin{column}{0.5\textwidth}
\begin{shadedbox}[title=\textbf{Investment Memo Skill}]
\begin{itemize}\small
  \item Standard structure: exec summary, financials, risks, recommendation
  \item Required data points and formatting
  \item Every memo follows the same template
\end{itemize}
\end{shadedbox}
\end{column}
\end{columns}
\vspace{0.3cm}
\alert{Every analyst produces comparable, auditable outputs.}
\end{frame}

\begin{frame}{CLAUDE.md: Global vs.\ Per-Project}
Claude reads \texttt{CLAUDE.md} files automatically at the start of every session --- no need to repeat instructions.
\vspace{0.3cm}
\begin{columns}[T]
\begin{column}{0.5\textwidth}
\begin{shadedbox}[title=\textbf{Global (\textasciitilde/.claude/CLAUDE.md)}]
\begin{itemize}\small
  \item Applies to \alert{every} project
  \item Output style preferences (chart colors, file formats)
  \item Preferred libraries (e.g., ``use pandas, not polars'')
  \item Personal conventions
\end{itemize}
\end{shadedbox}
\end{column}
\begin{column}{0.5\textwidth}
\begin{shadedbox}[title=\textbf{Per-Project (project root)}]
\begin{itemize}\small
  \item Applies only to \alert{this} folder
  \item Data context (column names, units, sources)
  \item Domain guidance (industry terms, formulas)
  \item Project-specific workflows
\end{itemize}
\end{shadedbox}
\end{column}
\end{columns}
\vspace{0.3cm}
Both are plain text files.  Claude reads the global file first, then the project file.  Project instructions can override global ones.
\end{frame}

\begin{frame}{Pre-Installed Skills: Document Creation}
\textbf{Skills} are packaged instructions that teach Claude Code a specific workflow.  The class installer pre-configured four.  Invoke any skill by typing \texttt{/} followed by its name.
\vspace{0.3cm}
\begin{columns}[T]
\begin{column}{0.32\textwidth}
\begin{shadedbox}[title=\textbf{/xlsx}]
\begin{itemize}\footnotesize
  \item Create and edit Excel workbooks
  \item Uses \alert{live formulas}, not hardcoded values
  \item Supports charts, formatting, multiple sheets
\end{itemize}
\end{shadedbox}
\end{column}
\begin{column}{0.32\textwidth}
\begin{shadedbox}[title=\textbf{/docx}]
\begin{itemize}\footnotesize
  \item Create and edit Word documents
  \item Professional formatting, headers, tables
  \item Supports tracked changes and comments
\end{itemize}
\end{shadedbox}
\end{column}
\begin{column}{0.32\textwidth}
\begin{shadedbox}[title=\textbf{/pptx}]
\begin{itemize}\footnotesize
  \item Create and edit PowerPoint decks
  \item Slide layouts, charts, speaker notes
  \item Can modify existing presentations
\end{itemize}
\end{shadedbox}
\end{column}
\end{columns}
\vspace{0.3cm}
\textbf{Example:} \textit{``/xlsx Create a loan amortization table for a \$200K mortgage at 6.5\% with monthly payments and a chart of principal vs.\ interest.''}
\end{frame}

\begin{frame}{Pre-Installed Skills: Skill Creator}
The \texttt{/skill-creator} skill is a guide for \alert{building your own skills}.
\vspace{0.3cm}
\begin{columns}[T]
\begin{column}{0.5\textwidth}
\begin{shadedbox}[title=\textbf{What It Does}]
\begin{itemize}\small
  \item Walks you through designing a new skill step by step
  \item Generates the skill definition file
  \item Installs the skill so it appears as a \texttt{/} command
\end{itemize}
\end{shadedbox}
\end{column}
\begin{column}{0.5\textwidth}
\begin{shadedbox}[title=\textbf{When to Use It}]
\begin{itemize}\small
  \item You repeat the same workflow across sessions
  \item You want Claude to follow a specific multi-step process
  \item You want to share a workflow with classmates
\end{itemize}
\end{shadedbox}
\end{column}
\end{columns}
\vspace{0.3cm}
We will use \texttt{/skill-creator} in the exercise below to build a personal task manager skill.
\end{frame}

\begin{frame}{Exercise: Create a Tasks Skill}
Use \texttt{/skill-creator} to build a \textbf{tasks} skill that manages personal and school to-do lists.
\vspace{0.3cm}
\begin{columns}[T]
\begin{column}{0.5\textwidth}
\begin{shadedbox}[title=\textbf{What the Skill Should Do}]
\begin{itemize}\small
  \item Store tasks in markdown files (e.g., \texttt{personal.md}, \texttt{school.md})
  \item Add new tasks with a category and due date
  \item List pending tasks when asked
  \item Mark tasks as complete
\end{itemize}
\end{shadedbox}
\end{column}
\begin{column}{0.5\textwidth}
\begin{shadedbox}[title=\textbf{Steps}]
\begin{enumerate}\small
  \item Type \texttt{/skill-creator} and describe the task manager you want
  \item Let Claude generate and install the skill
  \item Test it: ``/tasks Add: finish DCF exercise, due Tuesday''
  \item Test it: ``/tasks What's pending?''
\end{enumerate}
\end{shadedbox}
\end{column}
\end{columns}
\vspace{0.3cm}
This exercise teaches skill creation.  The same pattern applies to building skills for data retrieval, report generation, or any repeated workflow.
\end{frame}

\end{document}
